% Curs 10 - Pagina 6 (6-XII-1994).
% Fragment pentru % Curs 10 - Pagina 6 (6-XII-1994).
% Fragment pentru % Curs 10 - Pagina 6 (6-XII-1994).
% Fragment pentru % Curs 10 - Pagina 6 (6-XII-1994).
% Fragment pentru \input{courses/pages/curs10-pagina6-ro}.
% Fara preambul si fara intrerupere fortata de pagina.

\subsection*{Curs 10 - Pagina 6}

\noindent\underline{Def.}: Un grup $G$ s.n. \textbf{$p$-grup}, $p$ prim, dacă:
\[
(\forall)\,a\in G,\qquad o(a)=p^t,\quad t\geq 0.
\]
% In manuscris, cuvantul „p-grup” din prima linie apare taiat,
% desi este termenul definit de expresia care urmeaza.

\medskip
\noindent\underline{Obs.}: Dacă $G$ este grup finit, atunci
\[
\boxed{G\text{ este }p\text{-grup}\quad\Longleftrightarrow\quad
\operatorname{ord}(G)=p^n}.
\]

\noindent\underline{Dem.}: ``$\Leftarrow$'' evident.

\noindent ``$\Rightarrow$'': R.A. $\operatorname{ord}(G)\neq p^n$.
Atunci $(\exists)\,q$ prim, $q\neq p$, a.î. $q\mid\operatorname{ord}(G)$.

\noindent Conform teoremei lui Cauchy,
$(\exists)\,a\in G$ a.î. $o(a)=q$; cum $G$ este $p$-grup,
rezultă $p=q$, contradicţie.

\bigskip
\noindent-- Fie
\[
U_{p^n}=\{z\in\mathbb{C}\mid z^{p^n}=1\},
\qquad U_{p^n}\leq(\mathbb{C}^{*},\cdot).
\]
\noindent Şi $\operatorname{ord}(U_{p^n})=p^n$.

\[
\bigcup_{n\geq 1}U_{p^n}\ \leq\ (\mathbb{C}^{*},\cdot).
\]
\noindent Este $p$-grup infinit.

\bigskip
\noindent\underline{Def.}: Fie $G$ grup şi $H\leq G$.
$H$ e \underline{$p$-subgrup} al lui $G$ dacă şi numai dacă
$(H,\cdot)$ este $p$-grup.

\bigskip
\noindent\fbox{\textbf{T. lui Sylow 1}}.
Fie $G$ grup finit şi $\operatorname{ord}(G)=p^t m$,
$t\geq 1$ şi $p\nmid m$.

\medskip
\noindent Atunci $(\forall)\,i=\overline{1,t}$ există $H\leq G$
cu $|H|=p^i$.

\medskip
\noindent Mai mult, dacă $1\leq i<t$, atunci există $H\leq G$,
$|H|=p^i$, care este normal într-un subgrup al lui $G$
de ordin $p^{i+1}$.

\bigskip
\noindent\underline{Dem.}: Pentru $i=1$, O.K. din T. Cauchy.

\medskip
\noindent În aceleaşi ipoteze, un $p$-subgrup al lui $G$
de ordin $p^t$ s.n. \underline{$p$-subgrup Sylow} al lui $G$.
\quad (Not. $A=\operatorname{Syl}_p(G)$.)

\bigskip
\noindent\underline{Obs.}: Pe $A=\operatorname{Syl}_p(G)$,
$G\times A\longrightarrow A$ acţionează prin conjugare:
\[
P\in A,\ a\in G,\qquad
(a,P)\longmapsto aPa^{-1}
\quad\text{(conjugatul lui $P$ prin $a$)}.
\]
.
% Fara preambul si fara intrerupere fortata de pagina.

\subsection*{Curs 10 - Pagina 6}

\noindent\underline{Def.}: Un grup $G$ s.n. \textbf{$p$-grup}, $p$ prim, dacă:
\[
(\forall)\,a\in G,\qquad o(a)=p^t,\quad t\geq 0.
\]
% In manuscris, cuvantul „p-grup” din prima linie apare taiat,
% desi este termenul definit de expresia care urmeaza.

\medskip
\noindent\underline{Obs.}: Dacă $G$ este grup finit, atunci
\[
\boxed{G\text{ este }p\text{-grup}\quad\Longleftrightarrow\quad
\operatorname{ord}(G)=p^n}.
\]

\noindent\underline{Dem.}: ``$\Leftarrow$'' evident.

\noindent ``$\Rightarrow$'': R.A. $\operatorname{ord}(G)\neq p^n$.
Atunci $(\exists)\,q$ prim, $q\neq p$, a.î. $q\mid\operatorname{ord}(G)$.

\noindent Conform teoremei lui Cauchy,
$(\exists)\,a\in G$ a.î. $o(a)=q$; cum $G$ este $p$-grup,
rezultă $p=q$, contradicţie.

\bigskip
\noindent-- Fie
\[
U_{p^n}=\{z\in\mathbb{C}\mid z^{p^n}=1\},
\qquad U_{p^n}\leq(\mathbb{C}^{*},\cdot).
\]
\noindent Şi $\operatorname{ord}(U_{p^n})=p^n$.

\[
\bigcup_{n\geq 1}U_{p^n}\ \leq\ (\mathbb{C}^{*},\cdot).
\]
\noindent Este $p$-grup infinit.

\bigskip
\noindent\underline{Def.}: Fie $G$ grup şi $H\leq G$.
$H$ e \underline{$p$-subgrup} al lui $G$ dacă şi numai dacă
$(H,\cdot)$ este $p$-grup.

\bigskip
\noindent\fbox{\textbf{T. lui Sylow 1}}.
Fie $G$ grup finit şi $\operatorname{ord}(G)=p^t m$,
$t\geq 1$ şi $p\nmid m$.

\medskip
\noindent Atunci $(\forall)\,i=\overline{1,t}$ există $H\leq G$
cu $|H|=p^i$.

\medskip
\noindent Mai mult, dacă $1\leq i<t$, atunci există $H\leq G$,
$|H|=p^i$, care este normal într-un subgrup al lui $G$
de ordin $p^{i+1}$.

\bigskip
\noindent\underline{Dem.}: Pentru $i=1$, O.K. din T. Cauchy.

\medskip
\noindent În aceleaşi ipoteze, un $p$-subgrup al lui $G$
de ordin $p^t$ s.n. \underline{$p$-subgrup Sylow} al lui $G$.
\quad (Not. $A=\operatorname{Syl}_p(G)$.)

\bigskip
\noindent\underline{Obs.}: Pe $A=\operatorname{Syl}_p(G)$,
$G\times A\longrightarrow A$ acţionează prin conjugare:
\[
P\in A,\ a\in G,\qquad
(a,P)\longmapsto aPa^{-1}
\quad\text{(conjugatul lui $P$ prin $a$)}.
\]
.
% Fara preambul si fara intrerupere fortata de pagina.

\subsection*{Curs 10 - Pagina 6}

\noindent\underline{Def.}: Un grup $G$ s.n. \textbf{$p$-grup}, $p$ prim, dacă:
\[
(\forall)\,a\in G,\qquad o(a)=p^t,\quad t\geq 0.
\]
% In manuscris, cuvantul „p-grup” din prima linie apare taiat,
% desi este termenul definit de expresia care urmeaza.

\medskip
\noindent\underline{Obs.}: Dacă $G$ este grup finit, atunci
\[
\boxed{G\text{ este }p\text{-grup}\quad\Longleftrightarrow\quad
\operatorname{ord}(G)=p^n}.
\]

\noindent\underline{Dem.}: ``$\Leftarrow$'' evident.

\noindent ``$\Rightarrow$'': R.A. $\operatorname{ord}(G)\neq p^n$.
Atunci $(\exists)\,q$ prim, $q\neq p$, a.î. $q\mid\operatorname{ord}(G)$.

\noindent Conform teoremei lui Cauchy,
$(\exists)\,a\in G$ a.î. $o(a)=q$; cum $G$ este $p$-grup,
rezultă $p=q$, contradicţie.

\bigskip
\noindent-- Fie
\[
U_{p^n}=\{z\in\mathbb{C}\mid z^{p^n}=1\},
\qquad U_{p^n}\leq(\mathbb{C}^{*},\cdot).
\]
\noindent Şi $\operatorname{ord}(U_{p^n})=p^n$.

\[
\bigcup_{n\geq 1}U_{p^n}\ \leq\ (\mathbb{C}^{*},\cdot).
\]
\noindent Este $p$-grup infinit.

\bigskip
\noindent\underline{Def.}: Fie $G$ grup şi $H\leq G$.
$H$ e \underline{$p$-subgrup} al lui $G$ dacă şi numai dacă
$(H,\cdot)$ este $p$-grup.

\bigskip
\noindent\fbox{\textbf{T. lui Sylow 1}}.
Fie $G$ grup finit şi $\operatorname{ord}(G)=p^t m$,
$t\geq 1$ şi $p\nmid m$.

\medskip
\noindent Atunci $(\forall)\,i=\overline{1,t}$ există $H\leq G$
cu $|H|=p^i$.

\medskip
\noindent Mai mult, dacă $1\leq i<t$, atunci există $H\leq G$,
$|H|=p^i$, care este normal într-un subgrup al lui $G$
de ordin $p^{i+1}$.

\bigskip
\noindent\underline{Dem.}: Pentru $i=1$, O.K. din T. Cauchy.

\medskip
\noindent În aceleaşi ipoteze, un $p$-subgrup al lui $G$
de ordin $p^t$ s.n. \underline{$p$-subgrup Sylow} al lui $G$.
\quad (Not. $A=\operatorname{Syl}_p(G)$.)

\bigskip
\noindent\underline{Obs.}: Pe $A=\operatorname{Syl}_p(G)$,
$G\times A\longrightarrow A$ acţionează prin conjugare:
\[
P\in A,\ a\in G,\qquad
(a,P)\longmapsto aPa^{-1}
\quad\text{(conjugatul lui $P$ prin $a$)}.
\]
.
% Fara preambul si fara intrerupere fortata de pagina.

\subsection*{Curs 10 - Pagina 6}

\noindent\underline{Def.}: Un grup $G$ s.n. \textbf{$p$-grup}, $p$ prim, dacă:
\[
(\forall)\,a\in G,\qquad o(a)=p^t,\quad t\geq 0.
\]
% In manuscris, cuvantul „p-grup” din prima linie apare taiat,
% desi este termenul definit de expresia care urmeaza.

\medskip
\noindent\underline{Obs.}: Dacă $G$ este grup finit, atunci
\[
\boxed{G\text{ este }p\text{-grup}\quad\Longleftrightarrow\quad
\operatorname{ord}(G)=p^n}.
\]

\noindent\underline{Dem.}: ``$\Leftarrow$'' evident.

\noindent ``$\Rightarrow$'': R.A. $\operatorname{ord}(G)\neq p^n$.
Atunci $(\exists)\,q$ prim, $q\neq p$, a.î. $q\mid\operatorname{ord}(G)$.

\noindent Conform teoremei lui Cauchy,
$(\exists)\,a\in G$ a.î. $o(a)=q$; cum $G$ este $p$-grup,
rezultă $p=q$, contradicţie.

\bigskip
\noindent-- Fie
\[
U_{p^n}=\{z\in\mathbb{C}\mid z^{p^n}=1\},
\qquad U_{p^n}\leq(\mathbb{C}^{*},\cdot).
\]
\noindent Şi $\operatorname{ord}(U_{p^n})=p^n$.

\[
\bigcup_{n\geq 1}U_{p^n}\ \leq\ (\mathbb{C}^{*},\cdot).
\]
\noindent Este $p$-grup infinit.

\bigskip
\noindent\underline{Def.}: Fie $G$ grup şi $H\leq G$.
$H$ e \underline{$p$-subgrup} al lui $G$ dacă şi numai dacă
$(H,\cdot)$ este $p$-grup.

\bigskip
\noindent\fbox{\textbf{T. lui Sylow 1}}.
Fie $G$ grup finit şi $\operatorname{ord}(G)=p^t m$,
$t\geq 1$ şi $p\nmid m$.

\medskip
\noindent Atunci $(\forall)\,i=\overline{1,t}$ există $H\leq G$
cu $|H|=p^i$.

\medskip
\noindent Mai mult, dacă $1\leq i<t$, atunci există $H\leq G$,
$|H|=p^i$, care este normal într-un subgrup al lui $G$
de ordin $p^{i+1}$.

\bigskip
\noindent\underline{Dem.}: Pentru $i=1$, O.K. din T. Cauchy.

\medskip
\noindent În aceleaşi ipoteze, un $p$-subgrup al lui $G$
de ordin $p^t$ s.n. \underline{$p$-subgrup Sylow} al lui $G$.
\quad (Not. $A=\operatorname{Syl}_p(G)$.)

\bigskip
\noindent\underline{Obs.}: Pe $A=\operatorname{Syl}_p(G)$,
$G\times A\longrightarrow A$ acţionează prin conjugare:
\[
P\in A,\ a\in G,\qquad
(a,P)\longmapsto aPa^{-1}
\quad\text{(conjugatul lui $P$ prin $a$)}.
\]
